%% Research Statement — Dr. Panpan Cai
%% Updated: 2026-06-13.
%% This version expands the original 2019 NUS-Postdoc statement (still kept as
%% main_pre_edit_2026-06-13.tex) along two axes:
%%  (1) brings in the new home-service / mobile-manipulation domain following
%%      the "behaviour modelling -> real-time planning -> integrating learning"
%%      structure originally used for autonomous driving;
%%  (2) refreshes the autonomous-driving thread with Hi-Drive and Vec-QMDP,
%%      which carry POMDP planning into generic urban scenes (e.g. nuPlan).

\documentclass{article}
\usepackage[final]{nips_2018}
\usepackage[
  hidelinks,
  pdftitle={Research statement: Open-World Robot Planning and Learning},
  pdfauthor={Panpan Cai},
  pdfsubject={Research statement},
  pdfkeywords={robotics, planning under uncertainty, autonomous driving, mobile manipulation, foundation models},
]{hyperref}

\title{Open-World Robot Planning and Learning}

\author{%
  Panpan Cai \\
  Full-time Faculty, \href{https://www.sii.edu.cn/2025/1001/c82a512/page.htm}{Shanghai Innovation Institute} \\
  Associate Professor, \href{https://soai.sjtu.edu.cn/cn/facultydetails/zzjs/caipanpan}{School of Artificial Intelligence}, Shanghai Jiao Tong University \\
}

\begin{document}

\maketitle

My research focuses on \emph{large-scale decision making} for robots that must operate in \emph{complex, uncertain, and dynamic} real-world environments, so that they can interact seamlessly with humans and accomplish challenging long-horizon tasks. Decision making in such settings is fundamentally hard: robots have noisy sensors and limited models; the world is full of human-controlled agents with subtle and shifting intentions; and useful tasks span minute-long, multi-step plans whose state-action space is astronomical. Decisive progress requires combining the rigor of \emph{model-based planning} with the breadth of \emph{learning from large-scale data}.

Two application domains drive and stress-test the algorithms my group develops: \textbf{(i) autonomous driving in dense urban traffic}, where the robot must reason about many heterogeneous agents in real time, and \textbf{(ii) home-service mobile manipulation}, where the robot must interpret open-ended human instructions, reason over partially observed kitchens and apartments, and execute long-horizon plans built from contact-rich primitives. Concretely, my research advances three intertwined directions: \emph{modeling and reasoning under uncertainty}, \emph{real-time planning at scale}, and \emph{integrating planning with learning and foundation models}.

\subsection*{Modeling and reasoning under uncertainty}
A useful model not only captures the complexity of real-world dynamics and human behaviours, but also represents intrinsic uncertainties in principled ways. For \textit{autonomous driving}, I have developed traffic behavioural models that formalize interactions among heterogeneous agents as constrained optimization in the velocity space [1, 2], and an open-source simulator [3] that produces realistic driving simulations on real-world maps for end-to-end algorithm evaluation. For \textit{home-service robots}, the dominant uncertainties shift from continuous motion to symbolic and semantic ambiguity --- ambiguous human instructions, hidden or unknown object locations, and open-vocabulary object types. Tru-POMDP [4] tackles these challenges by combining a hierarchical \emph{Tree of Hypotheses}, queried from large language models, with an \emph{open-ended POMDP} formulation that supports rigorous Bayesian belief tracking and efficient belief-space planning. On complex object-rearrangement tasks across diverse kitchens, Tru-POMDP significantly outperforms LLM-based and LLM--tree-search hybrid planners, with stronger robustness to ambiguity and occlusion.

\subsection*{Real-time planning at scale}
Sophisticated planning brings combinatorial complexity --- the well-known ``curse of dimensionality'' and ``curse of history''. I attack this from a \emph{systems-and-algorithms co-design} angle. HyP-DESPOT [5, 6] introduces hybrid CPU/GPU parallel belief-tree search that speeds up large-scale POMDP planning by hundreds of times and is widely used as an open-source planner. Building on this foundation, my recent work pushes the \emph{practicalization} of POMDP planning on both target domains.

\textit{Generic urban driving.} Hi-Drive [7] formulates planning under behavioural and trajectory uncertainty as a hierarchical POMDP, treating driver models as high-level decision actions to manage exponential complexity, and refines trajectories via importance-sampling optimization over critical agents. On real-world urban benchmarks, Hi-Drive significantly outperforms state-of-the-art planning- and learning-based baselines across diverse driving situations. Vec-QMDP [8] further aligns POMDP search with modern CPUs' SIMD architecture via Data-Oriented Design and a hierarchical parallelism scheme (sub-trees across cores, vectorized expansion and collision checking within SIMD lanes), achieving 227x--1073x speedup over state-of-the-art serial planners and millisecond-level latency on large-scale benchmarks such as nuPlan. Together, Hi-Drive and Vec-QMDP move principled POMDP planning from constrained crowd-driving demos into the general urban driving stack, deployable on commodity CPUs without GPUs.

\textit{Mobile manipulation.} Neural Randomized Planning (NRP) [9] addresses long-range whole-body motion planning in cluttered household environments by combining global sampling-based motion planning with a local neural sampler. NRP exploits the search structure of the global planner to stitch together learned local sampling distributions into an adaptive global one, scaling to high-dimensional configuration spaces while preserving classical planning guarantees. Despite being trained only in simulation, NRP transfers zero-shot to a real robot operating in novel households.

\subsection*{Integrating planning with learning and foundation models}
Robots should improve from experience. Yet planning alone becomes intractable when problems are large-scale, while learning alone is brittle when domains shift. I have advocated a paradigm of \emph{``think locally, learn globally''} --- planning supplies online, problem-aware optimization, and learning amortizes global priors and heuristics.

\textit{Driving.} LeTS-Drive [10] learns global priors offline from tree search and uses them as heuristics to guide online belief-tree search, enabling sophisticated decision making in large, highly interactive crowds. LEADER [11] extends this idea by learning attention over driving behaviours for planning under uncertainty (CoRL 2022 Best Paper Finalist).

\textit{Home-service mobile manipulation.} The same paradigm has matured in the era of foundation models. UniDomain [12] pretrains a unified PDDL domain (3{,}137 operators, 2{,}875 predicates, 16{,}481 causal edges) from 12{,}393 manipulation videos; given a target task class, it retrieves and systematically fuses task-relevant atomics into compositional meta-domains. UniPlan [13] extends UniDomain to long-horizon mobile manipulation in large indoor environments by unifying scene topology, visuals, and robot capabilities into a holistic PDDL representation, grounded online by a vision-language model on a visual-topological scene map. AHAT [14] couples an LLM trained to map abstract human instructions plus textual scene graphs into PDDL subgoals with a new reinforcement-learning algorithm (TGPO) that injects external correction of intermediate reasoning into GRPO, yielding scalable long-horizon plans for human-style household tasks across very large environments.

\textit{Foundation models for physical intelligence.} At the policy level, my group studies how foundation models can be made physically capable. ForceVLA [15] treats external force sensing as a first-class modality in Vision-Language-Action models via a force-aware Mixture-of-Experts, improving contact-rich tasks (e.g.\ plug insertion) by 23.2\% over strong $\pi_0$-based baselines. MINT [16] argues that imitation learning should ``Mimic Intent, Not just Trajectories'': it disentangles behaviour intent from execution details via multi-scale frequency-space tokenization, and uses next-scale autoregression to enable one-shot skill transfer. I-Perceive [17] is a vision-language foundation model for active perception in mobile manipulation; it predicts task-relevant camera views from open-ended language instructions by fusing a VLM backbone with a geometric foundation model, generalizing zero-shot to novel scenes and tasks.

\section*{Summary and outlook}
I enjoy making principled approaches \emph{work} in the complex real world. The methodology my group has built --- belief-space planning, behaviour and intent modelling, massive parallelization, and tight coupling with learning --- was forged on autonomous driving and now also powers home-service mobile manipulation. The next phase of my research is to push the \emph{practicalization} of large-scale decision making one step further: open-world planners that handle truly unbounded objects, instructions, and scene types; foundation models that reason and plan in concert with classical solvers; and scalable systems that bring the rigor of model-based planning to settings where humans live and work alongside robots.

\section*{References}
\small

[1] Y. Luo, P. Cai, A. Bera, D. Hsu, W.S. Lee, D. Manocha. PORCA: Modeling and planning for autonomous driving among many pedestrians. \textit{IEEE RA-L}, 2018.

[2] Y. Luo*, P. Cai* (corresponding author), Y. Lee, D. Hsu. GAMMA: A General Agent Motion Model for Autonomous Driving. \textit{IEEE RA-L}, 2022.

[3] P. Cai*, Y. Lee*, Y. Luo, D. Hsu. SUMMIT: A Simulator for Urban Driving in Massive Mixed Traffic. \textit{IEEE ICRA}, 2020.

[4] W. Tang, X. He, Y. Huang, Y. Xiao, C. Lu, P. Cai (corresponding author). Tru-POMDP: Task Planning Under Uncertainty via Tree of Hypotheses and Open-Ended POMDPs. \textit{NeurIPS}, 2025.

[5] P. Cai, Y. Luo, D. Hsu, W.S. Lee. HyP-DESPOT: A Hybrid Parallel Algorithm for Online Planning under Uncertainty. \textit{RSS}, 2018.

[6] P. Cai, Y. Luo, D. Hsu, W.S. Lee. HyP-DESPOT: A Hybrid Parallel Algorithm for Online Planning under Uncertainty. \textit{IJRR}, 2021.

[7] X. Jin, C. Zeng, S. Zhu, C. Liu, P. Cai (corresponding author). Hi-Drive: Hierarchical POMDP Planning for Safe Autonomous Driving in Diverse Urban Environments. \textit{IEEE RA-L}, 2025.

[8] X. Jin, Y. Dong, B. Sun, H. Xu, Z. Hao, X. Lang, P. Cai (corresponding author). Vec-QMDP: Vectorized POMDP Planning on CPUs for Real-Time Autonomous Driving. \textit{RSS}, 2026.

[9] Y. Lu, Y. Ma, D. Hsu, P. Cai (corresponding author). Neural Randomized Planning for Whole Body Robot Motion. \textit{arXiv preprint}, 2024.

[10] P. Cai, Y. Luo, A. Saxena, D. Hsu, W.S. Lee. LeTS-Drive: Driving in a Crowd by Learning from Tree Search. \textit{RSS}, 2019.

[11] M.H. Danesh*, P. Cai* (corresponding author), D. Hsu. LEADER: Learning Attention over Driving Behaviors for Planning under Uncertainty. \textit{CoRL (Best Paper Finalist)}, 2022.

[12] H. Ye, Y. Xiao, C. Lu, P. Cai (corresponding author). UniDomain: Pretraining a Unified PDDL Domain from Real-World Demonstrations for Generalizable Robot Task Planning. \textit{NeurIPS}, 2025.

[13] H. Ye, Y. Xiao, C. Lu, P. Cai (corresponding author). UniPlan: Vision-Language Task Planning for Mobile Manipulation with Unified PDDL Formulation. \textit{arXiv preprint}, 2026.

[14] Z. Liu, Y. Li, R. Huang, C. Lu, P. Cai (corresponding author). Any House Any Task: Scalable Long-Horizon Planning for Abstract Human Tasks. \textit{arXiv preprint}, 2026.

[15] J. Yu*, H. Liu*, Q. Yu*, J. Ren, C. Hao, H. Ding, G. Huang, G. Huang, Y. Song, P. Cai, W. Zhang, C. Lu. ForceVLA: Enhancing VLA Models with a Force-aware MoE for Contact-rich Manipulation. \textit{NeurIPS}, 2025.

[16] R. Huang, C. Zeng, W. Tang, J. Cai, C. Lu, P. Cai (corresponding author). Mimic Intent, Not Just Trajectories. \textit{RSS}, 2026.

[17] Y. Huang, Z. Wang, W. Tang, C. Lu, P. Cai (corresponding author). I-Perceive: A Foundation Model for Vision-Language Active Perception. \textit{arXiv preprint}, 2026.

\end{document}
